\documentclass[12pt, letterpaper]{article}

\usepackage[utf8]{inputenc}
\usepackage[T1]{fontenc}
\usepackage[spanish]{babel}
\usepackage{geometry}
\usepackage{amsmath, amssymb, amsthm}
\usepackage{booktabs}
\usepackage{array}
\usepackage{microtype}
\usepackage{setspace}
\usepackage{parskip}
\usepackage{titlesec}
\usepackage{fancyhdr}
\usepackage{enumitem}
\usepackage{bm}
\usepackage{mathtools}
\usepackage{tcolorbox}
\usepackage{hyperref}

\tcbuselibrary{skins, breakable}

\geometry{top=2.8cm, bottom=3.0cm, left=3.2cm, right=3.2cm}
\setstretch{1.20}

\pagestyle{fancy}
\fancyhf{}
\fancyhead[C]{\small\textsc{Econometria · Gujarati --- Ejercicio 10.7}}
\fancyfoot[C]{\small\thepage}
\renewcommand{\headrulewidth}{0.4pt}
\renewcommand{\footrulewidth}{0pt}

\titleformat{\section}
  {\large\bfseries\scshape}{\thesection.}{0.6em}{}
  [\vspace{-4pt}\rule{\linewidth}{0.4pt}]
\titleformat{\subsection}{\normalsize\bfseries}{\thesubsection.}{0.5em}{}
\titleformat{\subsubsection}{\normalsize\itshape}{}{0em}{}

\newtheorem{prop}{Proposicion}
\newtheorem{lema}{Lema}
\theoremstyle{remark}
\newtheorem{obs}{Observacion}

\newtcolorbox{clave}{
  colback=white, colframe=black,
  boxrule=0.5pt, arc=3pt,
  left=8pt, right=8pt, top=6pt, bottom=6pt,
  breakable
}

\newcommand{\E}{\mathbb{E}}
\newcommand{\Var}{\operatorname{Var}}
\newcommand{\Cov}{\operatorname{Cov}}
\newcommand{\Corr}{\operatorname{Corr}}
\newcommand{\plim}{\operatorname{plim}}

\begin{document}

\begin{center}
  {\LARGE\bfseries Multicolinealidad en Series de Tiempo Economicas:}\\[6pt]
  {\LARGE\bfseries Tendencias, Inercia y Ciclos Comunes}\\[12pt]
  {\large\itshape Por que se sospecha multicolinealidad en PNB, oferta monetaria,}\\
  {\large\itshape precios, ingreso, desempleo y variables afines}\\[14pt]
  \rule{0.55\linewidth}{0.6pt}\\[8pt]
  {\normalsize Ejercicio~10.7 --- \textit{Econometria}, Gujarati \& Porter, 5.a ed.}\\[4pt]
  {\small\textsc{Analisis formal: tendencias estocasticas, autocorrelacion, factores comunes y diagnostico}}
\end{center}

\vspace{10pt}

\noindent\textbf{La pregunta.}\quad
En modelos de regresion con datos de series de tiempo macroeconomicas ---producto
nacional bruto (PNB), oferta monetaria, nivel de precios, ingreso nacional,
tasa de desempleo, poblacion, riqueza, entre otros--- se sospecha sistematicamente
la presencia de multicolinealidad. Este ejercicio justifica formalmente esa
sospecha identificando los mecanismos generadores de colinealidad propios de
las series temporales economicas.

\section{Naturaleza de las series de tiempo macroeconomicas}

\subsection{Definicion operativa de multicolinealidad en series de tiempo}

En el modelo de regresion multiple
\begin{equation}
  Y_t = \beta_1 + \beta_2 X_{2t} + \beta_3 X_{3t} + \cdots + \beta_k X_{kt} + u_t
  \label{eq:modelo}
\end{equation}
existe multicolinealidad cuando la matriz de correlaciones entre los regresores
$\{X_{2t}, \ldots, X_{kt}\}$ tiene elementos fuera de la diagonal cercanos
a~$\pm 1$. En terminos matriciales, la matriz $\mathbf{X'X}$ se aproxima a la
singularidad:
\begin{equation}
  \det(\mathbf{X'X}) \approx 0
  \label{eq:det}
\end{equation}
lo que infla los errores estandar de los estimadores MCO y hace que los
coeficientes individuales sean imprecisos o tengan signos incorrectos.

\subsection{Diferencia entre datos de corte transversal y series de tiempo}

En datos de corte transversal (observaciones de individuos, hogares o empresas
en un momento dado), las variables independientes exhiben variacion
\emph{idiosincratica}: el ingreso de un hogar no determina necesariamente la
riqueza del vecino. La correlacion entre regresores puede ser baja o moderada.

En series de tiempo macroeconomicas, en cambio, todas las variables se observan
sobre la \emph{misma} economia en \emph{momentos sucesivos}. Los mecanismos
macroeconomicos crean dependencia estructural entre las series. A continuacion
se formalizan las tres fuentes principales de esa dependencia.

\section{Fuente 1: tendencia temporal comun}

\subsection{Tendencia determinista}

Si varias variables macroeconomicas tienen tendencias deterministas de la misma
direccion, sus niveles estaran correlacionados aunque sus \emph{innovaciones}
sean independientes. Considere dos variables:
\begin{align}
  X_{2t} &= \alpha_2 + \delta_2 t + \varepsilon_{2t}
  \label{eq:trend2}\\
  X_{3t} &= \alpha_3 + \delta_3 t + \varepsilon_{3t}
  \label{eq:trend3}
\end{align}
con $\delta_2, \delta_3 > 0$ (ambas crecen en el tiempo) y
$\Cov(\varepsilon_{2t}, \varepsilon_{3t}) = 0$ (innovaciones independientes).
La correlacion entre $X_{2t}$ y $X_{3t}$ en niveles es:
\begin{equation}
  r_{23}
    = \frac{\Cov(X_{2t}, X_{3t})}{\sqrt{\Var(X_{2t})\Var(X_{3t})}}
    \approx \frac{\delta_2\delta_3 \Var(t)}
                 {\sqrt{(\delta_2^2\Var(t)+\sigma_2^2)(\delta_3^2\Var(t)+\sigma_3^2)}}
  \label{eq:corr_trend}
\end{equation}
Cuando $T$ es grande, $\Var(t) = (T^2-1)/12 \to \infty$ y el termino de
tendencia domina:
\begin{equation}
  r_{23} \xrightarrow{T\to\infty} \frac{\delta_2\delta_3}{|\delta_2||\delta_3|}
         = \text{signo}(\delta_2\delta_3) = +1
  \label{eq:corr_trend_lim}
\end{equation}
\textbf{Conclusion}: dos series con tendencia positiva tienen correlacion
asintoticamente igual a~1, aunque sus fluctuaciones sean independientes.
El PNB, el ingreso, la oferta monetaria y el nivel de precios crecen todos
en el largo plazo: su correlacion en niveles tiende a~1 por pura aritmetica
de las tendencias.

\subsection{Tendencia estocastica: caminata aleatoria}

Muchas series macroeconomicas son mejor descritas por procesos integrados
de orden~1 (I(1)), tambien llamados caminatas aleatorias con deriva:
\begin{align}
  X_{2t} &= \mu_2 + X_{2,t-1} + \varepsilon_{2t}
  \label{eq:rw2}\\
  X_{3t} &= \mu_3 + X_{3,t-1} + \varepsilon_{3t}
  \label{eq:rw3}
\end{align}
con $\varepsilon_{jt} \overset{\text{iid}}{\sim}(0, \sigma_j^2)$.
Expandiendo recursivamente desde $t=0$:
\begin{equation}
  X_{jt} = X_{j0} + \mu_j t + \sum_{s=1}^{t}\varepsilon_{js}
  \label{eq:rw_expand}
\end{equation}
La varianza crece con el tiempo: $\Var(X_{jt}) = \sigma_j^2 t \to \infty$.
En una muestra finita de $T$ periodos, dos series I(1) con derivas del mismo
signo ($\mu_2, \mu_3 > 0$) exhibiran una correlacion muestral alta aunque
$\Cov(\varepsilon_{2t}, \varepsilon_{3t}) = 0$.

Este fenomeno es la \textbf{regresion espuria} de Granger y Newbold (1974):
la regresion de $X_{2t}$ sobre $X_{3t}$ puede dar un $R^2$ alto y un
estadistico $t$ grande incluso cuando las dos series son generadas
independientemente, solo porque ambas tienen tendencias estocasticas.

\section{Fuente 2: autocorrelacion positiva e inercia}

\subsection{Suavidad de las series macroeconomicas}

Las series como el PNB, el consumo o la oferta monetaria no saltan
bruscamente de un trimestre al siguiente. Presentan \textbf{inercia}: el
valor pasado predice bien el valor presente. Formalmente, si cada variable
sigue un proceso AR(1):
\begin{equation}
  X_{jt} = \rho_j X_{j,t-1} + \varepsilon_{jt},
  \qquad |\rho_j| < 1
  \label{eq:AR1}
\end{equation}
la funcion de autocorrelacion (FAC) en el rezago $s$ es:
\begin{equation}
  \rho_{j,s} = \Corr(X_{jt}, X_{j,t-s}) = \rho_j^s
  \label{eq:FAC}
\end{equation}
Para datos trimestrales de variables nominales, $\rho_j \approx 0.9$--$0.95$.
Esto implica que $X_{jt}$ y $X_{j,t-4}$ tienen correlacion $\rho_j^4 \approx 0.66$.

\subsection{Por que la inercia genera colinealidad entre variables distintas}

Si dos variables $X_{2t}$ y $X_{3t}$ son ambas suaves (alta autocorrelacion)
y responden a los mismos shocks macroeconomicos (ciclo economico, politica
monetaria, demanda agregada), sus trayectorias se parecen:
\begin{equation}
  \Corr(X_{2t}, X_{3t})
    = \frac{\sum_{s=-\infty}^{\infty} a_s b_s \sigma_\eta^2}
           {\sqrt{\Var(X_{2t})\Var(X_{3t})}}
  \label{eq:corr_factor}
\end{equation}
donde $a_s$ y $b_s$ son los coeficientes de la representacion de media movil
de cada serie y $\eta_t$ es el shock comun. Cuanto mas dominante sea el factor
comun, mas alta sera la correlacion.

\subsection{Implicacion para la matriz $\mathbf{X'X}$}

Sea $\mathbf{x}_{jt} = X_{jt} - \bar X_j$ la desviacion respecto a la media.
El elemento $(j, k)$ de $\mathbf{X'X}$ (normalizado) es:
\begin{equation}
  \frac{\sum_t x_{jt} x_{kt}}{T} \approx \hat\gamma_{jk}(0)
  \label{eq:XpX_elem}
\end{equation}
la covarianza muestral contemporanea. Si $X_{jt}$ y $X_{kt}$ tienen
factor comun dominante, $\hat\gamma_{jk}(0)$ es grande relativo a
$\hat\gamma_{jj}(0)$ y $\hat\gamma_{kk}(0)$, y la correlacion se aproxima a~1.

\section{Fuente 3: ciclos comunes y factores macroeconomicos compartidos}

\subsection{El ciclo economico como generador de colinealidad}

Durante las fases expansivas del ciclo, el PNB, el empleo, el ingreso,
el consumo y la inversion crecen simultaneamente. Durante las contracciones,
todas caen juntas. Este co-movimiento es la definicion misma del ciclo
economico (Burns y Mitchell, 1946). En terminos de un modelo de factor comun:
\begin{align}
  X_{2t} &= \lambda_2 F_t + e_{2t}
  \label{eq:factor2}\\
  X_{3t} &= \lambda_3 F_t + e_{3t}
  \label{eq:factor3}
\end{align}
donde $F_t$ es el factor macroeconomico latente (ciclo), $\lambda_j$ son las
cargas del factor (\textit{factor loadings}) y $e_{jt}$ son los componentes
idiosincraticos con $\Cov(e_{2t}, e_{3t}) = 0$.

La correlacion entre $X_{2t}$ y $X_{3t}$ es:
\begin{equation}
  \Corr(X_{2t}, X_{3t})
    = \frac{\lambda_2 \lambda_3 \Var(F_t)}
           {\sqrt{(\lambda_2^2\Var(F_t)+\sigma_{e2}^2)(\lambda_3^2\Var(F_t)+\sigma_{e3}^2)}}
  \label{eq:corr_factor_expr}
\end{equation}
Si el factor comun domina ($\Var(F_t) \gg \sigma_{ej}^2$):
\begin{equation}
  \Corr(X_{2t}, X_{3t}) \approx \text{signo}(\lambda_2\lambda_3)
  \label{eq:corr_factor_lim}
\end{equation}
Cuando $\lambda_2$ y $\lambda_3$ tienen el mismo signo (ambas variables
son prociclicas, como PNB e ingreso), la correlacion es cercana a~1.

\subsection{Politica economica como fuente adicional}

Las variables de politica economica ---base monetaria, gasto publico, tasa
de interes de referencia--- suelen responder al estado del ciclo y a las
mismas variables macroeconomicas. Un banco central que sigue una regla de
Taylor ajusta la tasa de interes en funcion del PNB y la inflacion, que a su
vez estan correlacionadas entre si. Esto crea correlacion entre instrumentos
de politica y variables objetivo.

\section{Consecuencias econometricas formales}

\subsection{Inflacion de varianzas}

Bajo multicolinealidad con factor de inflacion $\text{VIF}_j$:
\begin{equation}
  \Var(\hat\beta_j) = \frac{\sigma^2}{S_{jj}} \cdot \text{VIF}_j,
  \qquad \text{VIF}_j = \frac{1}{1 - R_j^2}
  \label{eq:var_vif}
\end{equation}
donde $S_{jj} = \sum_t x_{jt}^2$ y $R_j^2$ es el $R^2$ de la regresion
auxiliar de $X_{jt}$ sobre los demas regresores. Con correlaciones entre
series de tiempo superiores a $0.95$, los VIF superan facilmente $20$--$100$.

\subsection{Estadisticos $t$ pequenos a pesar de $R^2$ alto}

El estadistico $t$ del $j$-esimo coeficiente es:
\begin{equation}
  t_j = \frac{\hat\beta_j}{\text{ee}(\hat\beta_j)}
       = \frac{\hat\beta_j}{\hat\sigma} \cdot \sqrt{\frac{S_{jj}}{\text{VIF}_j}}
  \label{eq:t_stat}
\end{equation}
Con VIF grande, $t_j$ se reduce aunque $\hat\beta_j$ sea economicamente grande.
El resultado es la paradoja clasica de series de tiempo: $R^2 > 0.95$ con
todos los $t_j$ no significativos.

\subsection{Inestabilidad de los coeficientes}

El estimador MCO puede escribirse como:
\begin{equation}
  \hat{\boldsymbol{\beta}} = \boldsymbol{\beta} + (\mathbf{X'X})^{-1}\mathbf{X'u}
  \label{eq:beta_error}
\end{equation}
La varianza de $\hat{\boldsymbol{\beta}}$ escala con $(\mathbf{X'X})^{-1}$.
Cuando $\det(\mathbf{X'X}) \approx 0$, los elementos de $(\mathbf{X'X})^{-1}$
son grandes: pequenas perturbaciones en los datos (agregar o quitar una
observacion, cambiar ligeramente la muestra) producen cambios grandes en
$\hat{\boldsymbol{\beta}}$. Esto explica la inestabilidad numerica bien
documentada en regresiones macroeconometricas de largo plazo.

\section{Diagnostico de multicolinealidad en series de tiempo}

\subsection{Indicadores practicos}

\begin{enumerate}[leftmargin=1.8em, itemsep=6pt]
  \item \textbf{Factor de inflacion de varianza (VIF):}
    \begin{equation}
      \text{VIF}_j = \frac{1}{1 - R_j^2} > 10 \quad \Rightarrow \quad
      \text{colinealidad severa}
      \label{eq:VIF_criterio}
    \end{equation}

  \item \textbf{Numero de condicion} de la matriz $\mathbf{X'X}$
    estandarizada:
    \begin{equation}
      \kappa = \sqrt{\frac{\lambda_{\max}}{\lambda_{\min}}} > 30
      \quad \Rightarrow \quad \text{colinealidad severa}
      \label{eq:kappa}
    \end{equation}

  \item \textbf{Indices de condicion de Belsley, Kuh y Welsch (1980):}
    valores superiores a~30 en alguno de los indices de condicion de las
    columnas de $\mathbf{X}$ indican problemas de colinealidad que afectan
    a los coeficientes correspondientes.

  \item \textbf{Correlacion entre series en niveles vs. en diferencias:}
    si la correlacion cae drasticamente al tomar primeras diferencias
    ($\Delta X_{jt} = X_{jt} - X_{j,t-1}$), la colinealidad original
    proviene de la tendencia comun, no de una relacion economica estructural.
\end{enumerate}

\subsection{Distincion crucial: colinealidad espuria vs. estructural}

\begin{itemize}[leftmargin=1.6em, itemsep=5pt]
  \item \textbf{Colinealidad espuria}: las series estan correlacionadas
    solo porque ambas tienen tendencia o son I(1). No hay una relacion
    economica causal entre ellas. En este caso, la correlacion desaparece
    en diferencias. El remedio es trabajar en diferencias o, si hay
    cointegracion, usar el modelo de correccion de error (MCE).

  \item \textbf{Colinealidad estructural}: las variables estan relacionadas
    por un mecanismo economico real (por ejemplo, el ingreso y la riqueza
    crecen juntos porque la riqueza acumulada es funcion del ingreso pasado).
    En este caso, la correlacion persiste en diferencias. El remedio requiere
    restricciones teoricas, datos adicionales o metodos como Koyck o Almon.
\end{itemize}

\section{Sintesis: por que se sospecha multicolinealidad en series macroeconomicas}

\begin{clave}
  La multicolinealidad es practicamente inevitable en regresiones con series
  de tiempo macroeconomicas por tres razones que actuan simultaneamente:

  \begin{enumerate}[leftmargin=1.4em, itemsep=5pt]
    \item \textbf{Tendencia comun}: el PNB, el ingreso, la oferta monetaria
      y el nivel de precios crecen todos en el largo plazo. Sus niveles estan
      correlacionados asintoticamente con~1 aunque sus innovaciones sean
      independientes. La correlacion no refleja causalidad: es un artefacto
      de la tendencia compartida.

    \item \textbf{Inercia y autocorrelacion}: las series macroeconomicas son
      suaves ($\rho \approx 0.9$). En muestras de tamano moderado (20--50
      anos de datos trimestrales), los rezagos de una variable y otras
      variables del mismo ciclo tienen correlaciones muestrales altas incluso
      si las correlaciones poblacionales son moderadas.

    \item \textbf{Factor macroeconomico comun}: el ciclo economico, la
      politica monetaria y los shocks de oferta mueven simultaneamente todas
      las variables agregadas en la misma direccion. Las cargas del factor
      comun ($\lambda_j$) son positivas para la mayoria de las variables
      prociclicas (PNB, empleo, ingreso, consumo), lo que produce correlaciones
      altas entre todos los pares.
  \end{enumerate}

  La consecuencia practica es que en un modelo con varias de estas variables
  como regresores, los estimadores MCO tendran errores estandar grandes,
  coeficientes potencialmente con signo incorrecto y una marcada sensibilidad
  a la especificacion exacta del modelo, incluso con $R^2$ cercano a~1.
\end{clave}

\appendix

\section{Regresion espuria: formalizacion de Granger-Newbold}

Sean $X_{2t}$ y $X_{3t}$ dos caminatas aleatorias independientes sin deriva:
\begin{align}
  X_{2t} &= X_{2,t-1} + \varepsilon_{2t}\\
  X_{3t} &= X_{3,t-1} + \varepsilon_{3t}
\end{align}
con $\varepsilon_{2t} \perp \varepsilon_{3t}$. Por construccion,
$\beta_3 = 0$ en la regresion $X_{2t} = \beta_1 + \beta_3 X_{3t} + u_t$.
Sin embargo, Granger y Newbold (1974) demostraron que el estadistico~$t$ de
$\hat\beta_3$ no converge a su distribucion nula $t_{T-2}$. En cambio:
\begin{equation}
  t_{\hat\beta_3} \xrightarrow{d} \frac{\int_0^1 W_2\, dW_3}
                                        {\left(\int_0^1 W_3^2\,ds\right)^{1/2}}
  \label{eq:t_espurio}
\end{equation}
donde $W_2$ y $W_3$ son movimientos brownianos estandar independientes.
Esta distribucion no es la $t$ de Student: tiene colas mucho mas pesadas
y rechaza $H_0: \beta_3 = 0$ con probabilidad mucho mayor que el nivel
nominal. El $R^2$ de la regresion espuria converge en probabilidad a~1
a medida que $T \to \infty$, lo que produce la apariencia de una relacion
perfecta entre variables completamente independientes.

\section{Consecuencia del numero de condicion sobre la precision numerica}

Sea $\kappa(\mathbf{X'X})$ el numero de condicion espectral. La precision
relativa de $\hat{\boldsymbol{\beta}}$ esta acotada por:
\begin{equation}
  \frac{\|\delta\hat{\boldsymbol{\beta}}\|}{\|\hat{\boldsymbol{\beta}}\|}
  \leq \kappa(\mathbf{X'X}) \cdot
  \frac{\|\delta\mathbf{X'y}\|}{\|\mathbf{X'y}\|}
  \label{eq:cond_bound}
\end{equation}
donde $\delta$ denota perturbaciones en los datos. Si $\kappa = 100$, un
error relativo del $1\%$ en $\mathbf{X'y}$ (por redondeo o por una
observacion atipica) puede inducir un error relativo del $100\%$ en
$\hat{\boldsymbol{\beta}}$. Para series de tiempo macroeconomicas con
correlaciones de $0.99$, $\kappa$ puede superar $1000$, haciendo los
estimadores numericamente inutiles.

\vspace{14pt}
\noindent\rule{\linewidth}{0.4pt}

\noindent{\small\textit{Nota.} El fenomeno de regresion espuria fue
documentado por Granger y Newbold (1974) y formalizado asintoticamente
por Phillips (1986). El tratamiento correcto de series I(1) requiere
contrastar cointegracion (Engle-Granger, 1987; Johansen, 1991) antes
de decidir si trabajar en niveles o en diferencias.}

\end{document}
